\scanpage{56}
\mspage{12}
\fbox{\deleted{\unsure{$4^\circ$)|A)}}} \inserted{(B)}
\uline{Operation} \uline{of} \deleted{\uline{weakening \unsure{and|or}} of} \inserted{\uline{over}}\uline{marking}

\deleted{It can be defined \inserted{in terms of} by giving, \unsure{a|a} \inserted{\unsure{a given system of|connected component of}} for each edge of the contour $R$ \unsure{$R\ne\phi$} $\widetilde\Sigma\setminus R$ Let $\deleted{\unsure{\widetilde\Sigma_*|\Sigma_*}}\;\unsure{A_0^+|A_*^+}\unsure{\subset|\in}\unsure{A_0|A_*}$ be the set}
\deleted{\unsure{of the} \unsure{a}}
Let $A^+$ be the set of ``actually marked'' edges \inserted{(free or bound)}, \inserted{\unsure{of the|of} bipartite $\Gamma$ such that $r(a)>0$},
\deleted{and $B$ the set of},
corresponding to $C\in\pi_0(\unsure{\widetilde\Sigma|\Sigma})$ such that $C\cap R\ne\phi$, let $B^+$ be the set of connected components \unsure{of the|of}
\[
 \Sigma^+\setminus R\cap\Sigma^+
\]
(which are open intervals), an overmarking is defined by giving a map
\[
 B^+\xrightarrow{r'}\mathbb N,
\]
and can be described \deleted{by} \inserted{as} the operation which on each
\uline{edge} \inserted{\unsure{(open)|\deleted{(open)}} \deleted{which}}
of $\Sigma$, corresponding to a $b\in B^+$, $r'(b)$ new \deleted{\unsure{marking}} \inserted{rigidification} vertices.
This time, \deleted{\unsure{the operation}} there are arbitrary choices --- nevertheless the operation gives an object \inserted{say $(X,\Sigma,\unsure{R_1|\widehat R_1|R^{\prime}_1})$} of $S\mathcal TD^+$, defined up to unique isomorphism --- equipped of course with a marked set $R'\subset R_1$, such that the inverse operation of erasing $R'$ gives back the starting object $(X,\Sigma,R)$.

\margin{\deleted{edges} \deleted{\unsure{marked}} \inserted{actually} marked circles, \inserted{\unsure{corr.|comp.}} \uline{edges} on $\Sigma$ \unsure{that}}
\margin{(\uline{NB} they \unsure{are the intersections of|lie on|are derived from} the actually marked circles, and are open intervals)}
\mspage{55}
